\documentclass[12pt]{article}
\usepackage[utf8]{inputenc}
\usepackage[T1]{fontenc}
\usepackage{hyperref}
\title{Gapapaiwa Language Resources for K-12 Curriculum Development}
\author{Compiled by Claude Sonnet 5.5 (claude-sonnet-5-5) for Andr\'as Kornai}
\date{}
\begin{document}
\maketitle

\section{Introduction}
Gapapaiwa (ISO 639-3 pwg; also Gapa, Paiwa, Manape) is an Austronesian language of the Western Oceanic branch, spoken in Milne Bay Province, Papua New Guinea, on the Cape Vogel peninsula (south coast villages and inland along the Ruaba River). Wikipedia gives about 3,000 speakers; Joshua Project gives 5,400 for the people group. The Gapa and Paiwa dialects are close to Ghayavi.

\section{Primary Educational Resources}
No primers, textbooks or school programmes were found in search results.

\section{Linguistic and Reference Materials}
\begin{itemize}
\item Wikipedia, ``Gapapaiwa language''. \url{https://en.wikipedia.org/wiki/Gapapaiwa_language}
\item Omniglot page on Gapapaiwa (classification and dialects). \url{https://www.omniglot.com/writing/gapapaiwa.htm}
\item Max Planck EVA numeral survey page for Gapapaiwa, reporting a quinary counting system (1997 survey). \url{https://lingweb.eva.mpg.de/channumerals/Gapapaiwa.htm}
\end{itemize}

\section{Religious and Community Resources}
\begin{itemize}
\item A New Testament exists in the Gapa dialect, available in PDF form at eBible: \url{https://ebible.org/pdf/pwg/pwg_all.pdf} and \url{https://ebible.org/pdf/pwg/pwg_a4.pdf}. These are the largest known continuous text in the language.
\item Joshua Project, Paiwa/Manape people group in Papua New Guinea. \url{https://joshuaproject.net/people_groups/14261}
\end{itemize}

\section{Digital and Audio Resources}
No audio or video resources were seen in search results.

\section{Curriculum Development Context}
The language is vigorous, and used by the whole community, but is not known to be taught in schools. The New Testament text could serve as a corpus for developing literacy materials, though its orthography and register should be checked with speakers. Material is otherwise sparse, so new documentation would be needed.
\end{document}